\chapter{이론적 배경 및 관련 연구}
\label{ch:background}

\section{이론적 배경}

이 장에서는 본 연구와 관련된 이론적 배경을 설명합니다.

\subsection{핵심 개념}

핵심 개념에 대한 설명을 기술합니다. 표와 그림은 다음과 같이 작성합니다.

\begin{table}[h]
  \centering
  \caption{샘플 표}
  \label{tab:sample}
  \begin{tabular}{lcc}
    \toprule
    항목 & 값1 & 값2 \\
    \midrule
    A & 1.23 & 4.56 \\
    B & 7.89 & 0.12 \\
    C & 3.45 & 6.78 \\
    \bottomrule
  \end{tabular}
\end{table}

표~\ref{tab:sample}과 같이 상호참조가 가능합니다.

\section{관련 연구}

관련 연구를 리뷰합니다~\cite{sample2026}.

\subsection{선행 연구 A}

\subsection{선행 연구 B}

\section{연구의 차별성}

본 연구가 기존 연구와 어떻게 차별화되는지 기술합니다.
